\documentclass[a4paper,12pt,twoside]{memoir}

% Castellano
\usepackage[spanish,es-tabla]{babel}
\selectlanguage{spanish}
\usepackage[utf8]{inputenc}
\usepackage[T1]{fontenc}
\usepackage{lmodern} % scalable font
\usepackage{microtype}
\usepackage{placeins}
\usepackage{lscape}

\usepackage[backend=bibtex,dateabbrev=false]{biblatex}
\addbibresource{bibliografiaAnexos.bib}

\RequirePackage{booktabs}
\RequirePackage[table]{xcolor}
\RequirePackage{xtab}
\RequirePackage{multirow}

% Links
\PassOptionsToPackage{hyphens}{url}\usepackage[colorlinks]{hyperref}
\hypersetup{
	allcolors = {red}
}

% Ecuaciones
\usepackage{amsmath}

% Rutas de fichero / paquete
\newcommand{\ruta}[1]{{\sffamily #1}}

% Párrafos
\nonzeroparskip

% Huérfanas y viudas
\widowpenalty100000
\clubpenalty100000

% Evitar solapes en el header
\nouppercaseheads

% Imagenes
\usepackage{graphicx}
\newcommand{\imagen}[2]{
	\begin{figure}[!h]
		\centering
		\includegraphics[width=0.9\textwidth]{#1}
		\caption{#2}\label{fig:#1}
	\end{figure}
	\FloatBarrier
}

\newcommand{\imagenflotante}[2]{
	\begin{figure}%[!h]
		\centering
		\includegraphics[width=0.9\textwidth]{#1}
		\caption{#2}\label{fig:#1}
	\end{figure}
}



% El comando \figura nos permite insertar figuras comodamente, y utilizando
% siempre el mismo formato. Los parametros son:
% 1 -> Porcentaje del ancho de página que ocupará la figura (de 0 a 1)
% 2 --> Fichero de la imagen
% 3 --> Texto a pie de imagen
% 4 --> Etiqueta (label) para referencias
% 5 --> Opciones que queramos pasarle al \includegraphics
% 6 --> Opciones de posicionamiento a pasarle a \begin{figure}
\newcommand{\figuraConPosicion}[6]{%
  \setlength{\anchoFloat}{#1\textwidth}%
  \addtolength{\anchoFloat}{-4\fboxsep}%
  \setlength{\anchoFigura}{\anchoFloat}%
  \begin{figure}[#6]
    \begin{center}%
      \Ovalbox{%
        \begin{minipage}{\anchoFloat}%
          \begin{center}%
            \includegraphics[width=\anchoFigura,#5]{#2}%
            \caption{#3}%
            \label{#4}%
          \end{center}%
        \end{minipage}
      }%
    \end{center}%
  \end{figure}%
}

%
% Comando para incluir imágenes en formato apaisado (sin marco).
\newcommand{\figuraApaisadaSinMarco}[5]{%
  \begin{figure}%
    \begin{center}%
    \includegraphics[angle=90,height=#1\textheight,#5]{#2}%
    \caption{#3}%
    \label{#4}%
    \end{center}%
  \end{figure}%
}
% Para las tablas
\newcommand{\otoprule}{\midrule [\heavyrulewidth]}
%
% Nuevo comando para tablas pequeñas (menos de una página).
\newcommand{\tablaSmall}[5]{%
 \begin{table}
  \begin{center}
   \rowcolors {2}{gray!35}{}
   \begin{tabular}{#2}
    \toprule
    #4
    \otoprule
    #5
    \bottomrule
   \end{tabular}
   \caption{#1}
   \label{tabla:#3}
  \end{center}
 \end{table}
}

%
%Para el float H de tablaSmallSinColores
\usepackage{float}

%
% Nuevo comando para tablas pequeñas (menos de una página).
\newcommand{\tablaSmallSinColores}[5]{%
 \begin{table}[H]
  \begin{center}
   \begin{tabular}{#2}
    \toprule
    #4
    \otoprule
    #5
    \bottomrule
   \end{tabular}
   \caption{#1}
   \label{tabla:#3}
  \end{center}
 \end{table}
}

\newcommand{\tablaApaisadaSmall}[5]{%
\begin{landscape}
  \begin{table}
   \begin{center}
    \rowcolors {2}{gray!35}{}
    \begin{tabular}{#2}
     \toprule
     #4
     \otoprule
     #5
     \bottomrule
    \end{tabular}
    \caption{#1}
    \label{tabla:#3}
   \end{center}
  \end{table}
\end{landscape}
}

%
% Nuevo comando para tablas grandes con cabecera y filas alternas coloreadas en gris.
\newcommand{\tabla}[6]{%
  \begin{center}
    \tablefirsthead{
      \toprule
      #5
      \otoprule
    }
    \tablehead{
      \multicolumn{#3}{l}{\small\sl continúa desde la página anterior}\\
      \toprule
      #5
      \otoprule
    }
    \tabletail{
      \hline
      \multicolumn{#3}{r}{\small\sl continúa en la página siguiente}\\
    }
    \tablelasttail{
      \hline
    }
    \bottomcaption{#1}
    \rowcolors {2}{gray!35}{}
    \begin{xtabular}{#2}
      #6
      \bottomrule
    \end{xtabular}
    \label{tabla:#4}
  \end{center}
}

%
% Nuevo comando para tablas grandes con cabecera.
\newcommand{\tablaSinColores}[6]{%
  \begin{center}
    \tablefirsthead{
      \toprule
      #5
      \otoprule
    }
    \tablehead{
      \multicolumn{#3}{l}{\small\sl continúa desde la página anterior}\\
      \toprule
      #5
      \otoprule
    }
    \tabletail{
      \hline
      \multicolumn{#3}{r}{\small\sl continúa en la página siguiente}\\
    }
    \tablelasttail{
      \hline
    }
    \bottomcaption{#1}
    \begin{xtabular}{#2}
      #6
      \bottomrule
    \end{xtabular}
    \label{tabla:#4}
  \end{center}
}

%
% Nuevo comando para tablas grandes sin cabecera.
\newcommand{\tablaSinCabecera}[5]{%
  \begin{center}
    \tablefirsthead{
      \toprule
    }
    \tablehead{
      \multicolumn{#3}{l}{\small continúa desde la página anterior}\\
      \hline
    }
    \tabletail{
      \hline
      \multicolumn{#3}{r}{\small continúa en la página siguiente}\\
    }
    \tablelasttail{
      \hline
    }
    \bottomcaption{#1}
  \begin{xtabular}{#2}
    #5
   \bottomrule
  \end{xtabular}
  \label{tabla:#4}
  \end{center}
}


\newcommand{\hubu}{\hline \rule{0pt}{2.7ex}}

\definecolor{cgoLight}{HTML}{EEEEEE}
\definecolor{cgoExtralight}{HTML}{FFFFFF}

%
% Nuevo comando para tablas grandes sin cabecera.
\newcommand{\tablaSinCabeceraConBandas}[5]{%
  \begin{center}
    \tablefirsthead{
      \toprule
    }
    \tablehead{
      \multicolumn{#3}{l}{\small\sl continúa desde la página anterior}\\
      \hline
    }
    \tabletail{
      \hline
      \multicolumn{#3}{r}{\small\sl continúa en la página siguiente}\\
    }
    \tablelasttail{
      \hline
    }
    \bottomcaption{#1}
    \rowcolors[]{1}{cgoExtralight}{cgoLight}

  \begin{xtabular}{#2}
    #5
   \bottomrule
  \end{xtabular}
  \label{tabla:#4}
  \end{center}
}




\graphicspath{ {./img/} }

% Capítulos
\chapterstyle{bianchi}
\newcommand{\capitulo}[2]{
	\setcounter{chapter}{#1}
	\setcounter{section}{0}
	\chapter*{#2}
	\addcontentsline{toc}{chapter}{#2}
	\markboth{#2}{#2}
}

% Apéndices
\renewcommand{\appendixname}{Apéndice}
\renewcommand*\cftappendixname{\appendixname}

\newcommand{\apendice}[1]{
	%\renewcommand{\thechapter}{A}
	\chapter{#1}
}

\renewcommand*\cftappendixname{\appendixname\ }

% Formato de portada
\makeatletter
\usepackage{xcolor}
\newcommand{\tutor}[1]{\def\@tutor{#1}}
\newcommand{\course}[1]{\def\@course{#1}}
\definecolor{cpardoBox}{HTML}{E6E6FF}
\def\maketitle{
  \null
  \thispagestyle{empty}
  % Cabecera ----------------
\noindent\includegraphics[width=\textwidth]{cabecera}\vspace{1cm}%
  \vfill
  % Título proyecto y escudo informática ----------------
  \colorbox{cpardoBox}{%
    \begin{minipage}{.8\textwidth}
      \vspace{.5cm}\Large
      \begin{center}
      \textbf{TFG del Grado en Ingeniería Informática}\vspace{.6cm}\\
      \textbf{\LARGE\@title{}}
      \end{center}
      \vspace{.2cm}
    \end{minipage}

  }%
  \hfill\begin{minipage}{.20\textwidth}
    \includegraphics[width=\textwidth]{escudoInfor}
  \end{minipage}
  \vfill
  % Datos de alumno, curso y tutores ------------------
  \begin{center}%
  {%
    \noindent\LARGE
    Presentado por \@author{}\\ 
    en Universidad de Burgos --- \@date{}\\
    Tutor: \@tutor{}\\
  }%
  \end{center}%
  \null
  \cleardoublepage
  }
\makeatother


% Datos de portada
\title{UBUAssistant\\Interacción por voz con la plataforma Moodle \\Documentación Técnica}
\author{Álvaro Delgado Pascual}
\tutor{D. Raúl Marticorena Sánchez}
\date{\today}

\begin{document}

\maketitle



\cleardoublepage



%%%%%%%%%%%%%%%%%%%%%%%%%%%%%%%%%%%%%%%%%%%%%%%%%%%%%%%%%%%%%%%%%%%%%%%%%%%%%%%%%%%%%%%%



\frontmatter


\clearpage

% Indices
\tableofcontents

\clearpage

\listoffigures

\clearpage

\listoftables

\clearpage

\mainmatter

\appendix

\apendice{Plan de Proyecto Software}

\section{Introducción}

\section{Planificación temporal}

\section{Estudio de viabilidad}

\subsection{Viabilidad económica}

\subsection{Viabilidad legal}



\apendice{Especificación de Requisitos}

\section{Introducción}

\section{Objetivos generales}

\section{Catalogo de requisitos}

\section{Especificación de requisitos}



\apendice{Especificación de diseño}

\section{Introducción}

\section{Diseño de datos}

Los datos están correlacionados con el formato con el que se guardan en Moodle, aunque solo están modelados aquellos datos que me son útiles. En el paquete /src/model/ están las clases que modelan estos datos y son:
\begin{itemize}
	\item \textbf{User} para guardar datos relativos al usuario.
	\item \textbf{Course} para guardar los datos de los cursos.
	\item \textbf{Forum} para los foros de un curso.
	\item \textbf{Discussion} para las discusiones de un foro.
	\item \textbf{GradeItem} para las notas.
	\item \textbf{Event} para los eventos del calendario.
\end{itemize}

\section{Diseño procedimental}



\section{Diseño arquitectónico}

El proyecto realmente consiste en dos aplicaciones, una la aplicación con la interfaz gráfica y la otra sería Mycroft. Ambas aplicaciones utilizan una arquitectura cliente-servidor, que es una arquitectura en la que el cliente realiza peticiones a través de un medio y el servidor envía una respuesta en función de la petición a través del mismo medio. La aplicación gráfica actúa como cliente y el servidor es la API REST de la plataforma de Moodle, es decir, los \textit{web services}. Mycroft tiene como cliente sus \textit{skills} que se comunican con dos servidores, uno para las funciones de STT y otro para TTS. Y entre las dos aplicaciones también existe una arquitectura cliente-servidor, siendo la aplicación gráfica el servidor y los clientes las diferentes \textit{skills}.
\apendice{Documentación técnica de programación}

\section{Introducción}

A continuación se va a detallar la estructura del proyecto, como montar el entorno de desarrollo, como ejecutar la aplicación y como realizar pruebas.

\section{Estructura de directorios}

El proyecto está dividido de esta forma:

\begin{itemize}
	\item /latex/: contiene los ficheros que se necesitan para generar le PDF con la memoria y los anexos, así como las imágenes utilizadas y la bibliografía.
	\item /src/: contiene los ficheros fuentes del proyecto..
	\item /src/GUI/: ficheros encargados del cliente gráfico
	\item /src/model/: ficheros para representar los datos de Moodle.
	\item /src/prototipo/: ficheros del prototipo para Alexa
	\item /src/webservice/: contiene los ficheros encargados de hacer las peticiones al servidor de Moodle para recuperar los datos.
	\item /src/util/: utilidades para los distintos módulos de la aplicación.
	\item /src/skills/: contiene las 3 \textit{Skills} del proyecto.	
\end{itemize}

\section{Manual del programador}

Para el entorno de desarrollo se necesita Python3 con sus dependencias, Mycroft y un editor de código fuente.

\subsection{Python 3}

Python 3 viene instalado por defecto en Ubuntu, pero necesitamos instalar las dependencias para la aplicación. Estas dependencias son PyQt5 y MycroftMessageBus y se instalan mediante los comandos:
\begin{itemize}
	\item \textbf{sudo apt-get install python3-pyqt5}
	\item \textbf{pip3 install mycroft-messagebus-client}
\end{itemize}

\subsection{Mycroft}
Antes de nada, necesitamos tener git instalado, se instala mediante el comando \textbf{sudo apt install git}, y necesitamos una cuenta de Mycroft, para ello ir a \href{https://mycroft.ai/}{su página} y crearla.

\imagen{crear_mycroft1}{Creación de cuenta de Mycroft} \imagen{crear_mycroft2}{Creación de cuenta de Mycroft}

Para instalar Mycroft hay que seguir los pasos en la sección \textbf{Getting Started} de la guía que hay en \href{https://github.com/MycroftAI/mycroft-core}{su repositorio}, que son:
\begin{enumerate}
	\item \textbf{cd \detokenize{~}}
	\item \textbf{git clone https://github.com/MycroftAI/mycroft-core.git}
	\item \textbf{cd mycroft-core}
	\item \textbf{bash dev\_setup.sh}
\end{enumerate}

Tras introducir el comando \textbf{bash dev\_setup.sh} se nos abrirá el proceso de instalación interactiva. Le vamos a dar a todo a \textit{Yes} (escribir la letra Y en la consola) 

Una vez instalado, continuar a la sección \textbf{Running Mycroft} de \href{https://github.com/MycroftAI/mycroft-core}{su repositorio} para ejecutarlo, que se hace mediante el los comandos:
\begin{enumerate}
	\item \textbf{cd \detokenize{~}/mycroft-core}
	\item \textbf{./start-mycroft.sh debug}
\end{enumerate}

Tras iniciar Mycroft tenemos que decir por voz cualquier cosa para que se inicie el proceso de enlazamiento. La aplicación nos dirá un código (si se ha iniciado con la opción de \textit{debug} también aparecerá en pantalla). Este código hay que ponerlo al añadir un nuevo dispositivo. En \href{https://home.mycroft.ai}{la página de Mycroft} en la esquina superior derecha, en la sección de \textbf{Devices}, \textbf{Add Device} en el apartado de \textbf{Pairing code}. Al añadir el nuevo dispositivo es importante que en el apartado \textbf{Voice} se seleccione \textbf{Google Voice}.

Para parar Mycroft si se ha iniciado con la opción de \textit{debug} se puede presionar la combinación de teclas \textbf{ctrl + c} o escribir en la interfaz \textbf{:exit}

\subsection{Editor}

Se puede utilizar el editor de código fuente con el que más cómodo se sienta cada uno, en esta sección voy a explicar como configurar el editor Atom que es el que he estado usando y el que más cómodo me parece, aunque hay alternativas como VisualStudio Code que son muy similares.

Para instalar Atom nos vamos a \href{https://atom.io/}{su página}, lo descargamos y lo ejecutamos. Una vez instalado vamos a instalar dos paquetes para que sea mucho más cómodo programar con Python. En la pestaña \textbf{Edit} hacemos click en \textbf{Preferences}.

\imagen{atom_preferences}{Preferencias Atom}

En esta nueva ventana que se nos ha abierto hacemos click en \textbf{Install} y en el buscador escribimos \textbf{atom-ide-ui} e instalamos el paquete. Hacemos lo mismo buscando \textbf{ide-python}.

\imagen{install_atom_ide}{Instalar atom-ide} \imagen{install_ide_python}{Instalar ide-python}

Para que funcione el paquete \textbf{ide-python} tenemos que instalar \textbf{python-language-server} escribiendo en la consola el comando:
\begin{itemize}
	\item \textbf{python3 -m pip install python-language-server[all]}
\end{itemize}

\imagen{pyls}{Instalar Python Language Server}

Ahora vamos a configurar el paquete, para ello hacemos click en la pestaña \textbf{Packages}, después en \textbf{Settings} de \textbf{ide-python}. En la nueva ventana que se nos ha abierto se puede configurar el comportamiento de todos los componentes de \textbf{python-language-server}, pero para que funcione de momento solo nos interesa cambiar la primera línea en la que pone python y poner \textbf{python3}

\imagen{ide_python_settings}{Configuración del paquete ide-python} \imagen{ide_python_settings2}{Configuración del paquete ide-python}


\section{Compilación, instalación y ejecución del proyecto}

Clonamos el repositorio en nuestra carpeta personal mediante los comandos:
\begin{enumerate}
	\item \textbf{cd \detokenize{~}}
	\item \textbf{git clone https://github.com/adp1002/UBUAssistant}
\end{enumerate}

Antes de nada hay que cambiar los ficheros \_\_init\_\_.py de las 3 skills (ubu-course, ubu-grades y ubu-course en la carpeta /src/skills/). En la línea 4, sys.path.append(expanduser('~') + '/UBUAssistant-1.X/src'), hay que cambiar /UBUAssistant-1.X/src a \textbf{/UBUAssistant/src}

Cuando Lo hayamos cambiado, mover el contenido de la carpeta \textbf{\detokenize{~}/src/skills/} a \textbf{/opt/mycroft/skills/}, ejecutando el comando:
\textbf{sudo mv \detokenize{~}/UBUAssistant-1.X/src/skills/* /opt/mycroft/skills/})

Para iniciar la aplicación vamos al directorio /src/ y ejecutamos mediante los comandos:
\begin{enumerate}
	\item \textbf{cd \detokenize{~}/UBUAssistant/src/}
	\item \textbf{python3 -m GUI.main >\null> logs.log}
\end{enumerate}

\section{Pruebas del sistema}

Para probar nueva funcionalidad introducida se puede utilizar la plataforma de pruebas de Moodle \href{https://school.moodledemo.net/}{Mount Orage School} o la plataforma con la que se vaya a utilizar la aplicación, que en mi caso es UBUVirtual y es con lo que he realizado pruebas principalmente.
\apendice{Documentación de usuario}

\section{Introducción}

En este apartado se explica como puede un usuario prepara el entorno para ejecutar la aplicación y cómo utilizar la aplicación.

\section{Requisitos de usuarios}

Sistema operativo GNU/Linux
Python3
Conexión a Internet

\section{Instalación}

\subsection{Ubuntu}

Si ya tienes Ubuntu u otra distribución Linux vete al siguiente apartado.

Para utilizar Ubuntu se va a necesitar algún software de virtualización. En esta guía se va a utilizar VirtualBox, que se puede descargar desde \href{https://www.virtualbox.org/wiki/Downloads}{su página}.
\imagen{descarga_virtualbox}{Descarga de VirtualBox}
Desde su página elegimos la opción de \textbf{Windows hosts} (u otra en función del sistema operativo que se esté utilizando). Esperamos a que se inicie la descarga y ejecutamos el archivo .exe que hemos descargado. Se nos abrirá el instalador y le damos a \textbf{Next, Next, Next, Yes, Install}.

Una vez instalado VirtualBox vamos a descargar Ubuntu. Vamos a \href{https://releases.ubuntu.com/18.04.4/}{su página} y hacemos click en el enlace que pone \textbf{64-bit PC (AMD64) desktop image}. Una vez haya finalizado la descarga abrimos VirtualBox. En la ventana de VirtualBox, hacemos click en \textbf{Nueva} y la configuramos para \textbf{Ubuntu (64-bit)}. A continuación elegimos la cantidad de memoria y creamos un disco virtual de, al menos, 30 GB.

\imagen{VM_crear0}{Nueva máquina virtual} \imagen{VM_crear1}{Tipo de sistema operativo}

Después hacemos click en el apartado de configuración en la ventana inicial de VirtualBox. En la pestaña \textbf{Almacenamiento} hacemos click en \textbf{Añadir una unidad óptica}, que es el botón a la derecha de \textbf{Controlador: IDE}. En la nueva ventana hacemos click en \textbf{Añadir} y seleccionamos el archivo .iso de Ubuntu. En la ventana de configuración también se puede ajustar el número de procesadores y la cantidad de memoria gráfica desde las pestañas \textbf{Sistema} y \textbf{Pantalla} respectivamente. Recomiendo utilizar, como mínimo, 2048MB de memoria RAM, 2 núcleos de procesador y 128MB de memoria de vídeo.

\imagen{VM_config0}{Configuración de máquina virtual} \imagen{VM_config1}{Parámetros de configuración}

Una vez terminada la configuración, iniciamos Ubuntu desde la ventana principal de VirtualBox con el botón \textbf{Iniciar}. En la primera pantalla de Ubuntu seleccionamos el idioma (español) y le damos a instalar Ubuntu.
Después seleccionamos el idioma del teclado y la instalación de Ubuntu, en nuestro caso es suficiente con la \textbf{instalación mínima}. A continuación borramos el disco virtual, seleccionamos la franja horaria y completamos los campos que nos pide. Una vez termine de instalar nos pedirá reiniciar y después darle al \textbf{Enter}. Cuando se haya iniciado Ubuntu solo queda instalar las \textbf{guest additions}, que es totalmente opcional pero aumenta el rendimiento. Para instalarlas hacer click en la pestaña \textbf{Dispositivos} y \textbf{Insertar imagen de CD de las <<Guest Additions>>...} y reiniciar el sistema operativo.

\imagen{Ubuntu_install}{Instalación de Ubuntu} \imagen{Ubuntu_install2}{Instalación de Ubuntu}

\subsection{Python3 y sus dependencias}

Python3 viene instalado por defecto en Ubuntu 18.04. Para instalar las dependencias necesitamos pip3, que se instala con el comando \textbf{sudo apt install python3-pip}. Ahora las dependencias que necesitamos se instalan mediante los comandos \textbf{sudo apt-get install python3-pyqt5} y \textbf{pip3 install mycroft-messagebus-client}

\subsection{Mycroft}
Antes de nada, necesitamos tener git instalado, se instala mediante el comando \textbf{sudo apt install git}, y necesitamos una cuenta de Mycroft, para ello ir a \href{https://home.mycroft.ai/}{su página} y crearla.

Para instalar Mycroft hay que seguir los pasos en la sección \textbf{Getting Started} de la guía que hay en \href{https://github.com/MycroftAI/mycroft-core}{su repositorio}, que son:
\begin{enumerate}
	\item \textbf{cd \detokenize{~}}
	\item \textbf{git clone https://github.com/MycroftAI/mycroft-core.git}
	\item \textbf{cd mycroft-core}
	\item \textbf{bash dev\_setup.sh}
\end{enumerate}

Tras introducir el comando \textbf{bash dev\_setup.sh} se nos abrirá el proceso de instalación interactiva. Le vamos a dar a todo a \textit{Yes} (escribir la letra Y en la consola) 

Una vez instalado, continuar a la sección \textbf{Running Mycroft} de \href{https://github.com/MycroftAI/mycroft-core}{su repositorio} para ejecutarlo, mediante el los comandos:
\begin{enumerate}
	\item \textbf{cd \detokenize{~}/mycroft-core}
	\item \textbf{./start-mycroft.sh debug}
\end{enumerate}

Tras iniciar Mycroft tenemos que decir por voz cualquier cosa para que se inicie el proceso de enlazamiento. La aplicación nos dirá un código (si se ha iniciado con la opción de \textit{debug} también aparecerá en pantalla). Este código hay que ponerlo al añadir un nuevo dispositivo. En \href{https://home.mycroft.ai}{la página de Mycroft} en la esquina superior derecha, en la sección de \textbf{Devices}, \textbf{Add Device} en el apartado de \textbf{Pairing code}. Al añadir el nuevo dispositivo es importante que en el apartado \textbf{Voice} se seleccione \textbf{Google Voice}.

Para parar Mycroft si se ha iniciado con la opción de \textit{debug} se puede presionar la combinación de teclas \textbf{ctrl + c} o escribir en la interfaz \textbf{:exit}

\imagen{cuenta_mycroft}{Crear dispositivo en Mycroft}

\section{Manual del usuario}

%En Ubuntu, para ejecutar la aplicación hay que descargar el \href{https://github.com/adp1002/UBUAssistant/releases}{proyecto de GitHub} (el archivo \textbf{Source code (tar.gz)}) y  descomprimirlo en la carpeta personal.

%\imagen{release}{Descarga de la release del proyecto}

Clonamos el repositorio en nuestra carpeta personal mediante los comandos:
\begin{enumerate}
	\item \textbf{cd \detokenize{~}}
	\item \textbf{git clone https://github.com/adp1002/UBUAssistant}
\end{enumerate}

Cuando tengamos el proyecto descargado, mover el contenido de la carpeta \textbf{\detokenize{~}/src/skills/} a \textbf{/opt/mycroft/skills/}, ejecutando el comando:
\textbf{sudo mv \detokenize{~}/UBUAssistant/src/skills/* /opt/mycroft/skills/})

Nos vamos a la carpeta \textbf{/src/} del proyecto y ejecutamos la aplicación mediante los comandos:
\begin{enumerate}
	\item \textbf{cd \detokenize{~}/UBUAssistant/src}
	\item \textbf{python3 -m ./GUI/main.py \detokenize{>}> logs.log \&}
\end{enumerate}
Se nos abrirá la ventana de inicio de sesión. En esta ventana se introducirán los credenciales del usuario para la plataforma de Moodle a la que se quiera conectar, así como la URL de esta plataforma.
Es importante que en el campo \textbf{host} se introduzca la URL exactamente con el formato que viene indicado. Por ejemplo, para conectarse a UBUVirtual la URL que habrá que introducir es, literalmente, \textbf{https://ubuvirtual.ubu.es}.
También se puede cambiar el idioma de la aplicación en la esquina superior derecha y recordar el usuario y host para futuras ejecuciones.

\imagen{login_window}{Ventana de inicio de sesión}

En la parte superior de la ventana principal de la aplicación se muestra la estructura de todos los \textit{intents}. El botón del micrófono sirve para alternar entre activar y desactivarlo, cuando el botón está verde está activo y cuando está rojo está inactivo.
Para usar la aplicación mediante voz, hay que decir la \textit{wake word} (por defecto es "Hey Mycroft") y se escuchará un sonido si ha detectado la \textit{wake word}. Después de escuchar el sonido se puede preguntar lo que queramos con la frase correspondiente. Para utilizar la conversación mediante texto hay que escribir la pregunta, sin la \textit{wake word}, y darle a enviar. Toda la conversación, ya sea mediante texto o voz, será transcrita en la ventana.

\imagen{main_window}{Ventana principal} \imagen{ejemplo_skill}{Ejemplo de salida de la aplicación}

Se pueden tener activas las \textit{skills} que queramos, haciendo click en el botón de \textbf{Administrar Skills}. Las \textit{skills} que están marcadas son las activas y las no marcadas no están activas. Para activar o desactivar una \textit{skill}, marcar o desmarcar las \textit{skills} que queramos y pulsar el botón \textbf{Guardar}

\imagen{skills_window}{Ventana de administración de las \textit{Skills}}

Con el botón de \textbf{Abrir Logs} podemos ver los distintos mensajes de información que se van guardando durante la ejecución de la aplicación. Los títulos de las pestañas hacen referencia al fichero de log que se muestra en esa pestaña.

\imagen{log_window}{Visualización de logs}

\printbibliography

\end{document}
